
@@M0@@ de @@M1@@-modules localement libres de rang @@M2@@) par recollement du @@M3@@-module libre @@M4@@ sur des cartes locales trivialisantes moyennant la représentation @@M5@@. De façon précise, à la suite @@M6@@ correspond un @@M7@@-torseur à droite @@M8@@ tel que l'on a un isomorphisme canonique
@@M9@@
où @@M10@@ désigne le produit contracté, et @@M11@@ opère sur @@M12@@ via la @@M13@@-ième projection @@M14@@ ; alors toute représentation
@@M15@@
fait opérer @@M16@@ sur @@M17@@ par des @@M18@@-automorphismes, et définit ainsi un multifoncteur @@M19@@ tel que
@@M20@@
De plus, pour toute représentation composée @@M21@@ on trouve
@@M22@@

\subsection*{6.2.}
Pour tout @@M23@@-module localement libre @@M24@@ de rang @@M25@@ nous désignerons par @@M26@@ le @@M27@@-module inversible @@M28@@. Gardant les notations de 6.1, le foncteur @@M29@@ est déterminé, à isomorphisme canonique près, par la représentation correspondante
@@M30@@

Supposons maintenant que @@M31@@ soit le topos annelé des faisceaux sur le site annelé @@M32@@ muni de la topologie étale. Alors on peut calculer @@M33@@, ou ce qui revient au même, la représentation @@M34@@, à l'aide de la

